Aufgrund von
[[Stetige Abbildung/Bild eines kompakten Raumes/Kompakt/Fakt|Fakt]]
ist

\mavergleichskette
{\vergleichskette
{ f(X)
}
{ \subseteq }{ \R
}
{  }{
}
{    }{
}
{   }{
}
}
{}{}{}
kompakt, also nach
[[Kompaktheit/Satz von Heine-Borel/Fakt|Fakt]]
\definitionsverweis {abgeschlossen}{}{}
und
\definitionsverweis {beschränkt}{}{.}
Insbesondere ist

\mavergleichskette
{\vergleichskette
{ f(X)
}
{ \leq }{ M
}
{  }{
}
{    }{
}
{   }{
}
}
{}{}{}
für eine reelle Zahl $M$. Wegen

\mavergleichskette
{\vergleichskette
{ X
}
{ \neq }{ \emptyset
}
{  }{
}
{    }{
}
{   }{
}
}
{}{}{}
besitzt $f(X)$ wegen
[[Reelle Zahlen/Beschränkte Teilmenge hat Supremum/Fakt|Fakt]]
ein
\definitionsverweis {Supremum}{}{}
$s$ in $\R$, das wegen der Abgeschlossenheit
nach [[Reelle Zahlen/Abgeschlossene nach oben beschränkte Menge/Enthält Supremum/Fakt|Fakt]]
zu $f(X)$ gehört, also das Maximum von $f(X)$ ist. Daher gibt es auch ein

\mavergleichskette
{\vergleichskette
{ x
}
{ \in }{ X
}
{  }{
}
{    }{
}
{   }{
}
}
{}{}{}
mit

\mavergleichskette
{\vergleichskette
{ f(x)
}
{ = }{ s
}
{  }{
}
{    }{
}
{   }{
}
}
{}{}{.}

 [[Kategorie:Latexseite]]
